% sec_fullpos.tex -- Supplemental Material section: full positivity of Omega^EE at finite mass.
% Meant to be \input into supplement.tex just before the section on the angular route (sec:angular).
\section{Full positivity of the energy--energy action at finite mass}
\label{sec:fullpos}

This section reports what is established about assumption (iii) for $\Omega^{EE}$ at finite $(\Delta,J)$, beyond the leading-Regge
floor T0 (Sec.~\ref{sec:T0}). It uses an exact rewriting of the heavy action in the detector-ladder variables of
Sec.~\ref{sec:angular}. Like Secs.~\ref{sec:pos}--\ref{sec:allgap}, it is evidence for (iii) and is never used as an input. Unless
stated otherwise $\Delta_g=1000$, and the $k=2$ weight is the published, endpoint-repaired continuum weight that defines
$\Omega^{EE}$ (Table~\ref{tab:data}). \emph{Summary.} Positivity of the finite-mass action is certified for every spin $J\le766$ at
every mass, and for every larger spin at impact parameter $\hat b\le400$ (Table~\ref{tab:fpcover}). At very large impact parameter a
contour argument gives it at grade [D-sketch]. The rest, large spin near the twist edge at moderate impact parameter, is open; numerical
evidence there is listed in Sec.~\ref{sec:fullpos}.E.

\emph{Notation.} $T:=T_\Delta=(\Delta-1)(\Delta-2)$ and $C:=J(J+1)$, so that $m^2=T-C$; $r:=C/T$, $w:=1-r=m^2/T$ and
$\ell:=-\ln(1-r)=\ln(1+C/m^2)$; $a:=\Delta-J-2$. The impact parameter is $\hat b:=(2J+1)\Lambda/m$, the $b$ of Sec.~\ref{sec:defs}.A,
hatted here because $b\in\{0,4\}$ labels the two helicity ladders of Sec.~\ref{sec:angular}. Then
$\ell=\ln\{1+(\hat b^2/4\Lambda^2)[1-(2J+1)^{-2}]\}$, so $\hat b=400$ is $\ell\simeq0.039$ and $\ell=10$ is $\hat b\simeq3.0\times10^5$.
On ladder $b$: $u_j:=(b+j)(b+j+1)$, $n:=J-b$, $D_k:=\binom{k+3}{3}$, $\sigma_0=+1$, $\sigma_4=-1$, and
\begin{equation}
U_{b,k}:=\prod_{j<k}\frac{C-u_j}{T-u_j}=\frac{(b-J)_k\,(J+b+1)_k}{(b+2-\Delta)_k\,(\Delta+b-1)_k},
\label{eq:fp-U}
\end{equation}
which vanishes for $k>n$ (the ladder terminates). For heavy states $0\le U_{b,k}\le r^k$.

\subsection{The target}
For a heavy primary $\mathcal O$ the action of $\Omega^{EE}$ is
\begin{equation}
\Omega^{EE}_{\mathcal O}=\mathcal B(\Delta,J)\sum_{b=0,4}|a_{\mathcal O,b,0}|^2\,H_{{\rm ang},b}(\Delta,J),\qquad
\mathcal B=\frac{2\pi G\Lambda^4}{9}\,\frac{T^4}{(T-C)^3}>0,
\label{eq:fp-target}
\end{equation}
with $a_{\mathcal O,b,0}$ the rest-frame detector amplitudes of Sec.~\ref{sec:angular} and $H_{{\rm ang},b}$ given in
Eq.~\eqref{eq:fp-H} below. There are no $b=0/4$ cross terms: the ladders have disjoint supports, the first-moment matrix is
tridiagonal, and the $k=3$ helicity-flip terms cancel on the physical ladders. In the basis of OPE coefficients the matrix
$M_\Omega(\mathcal O)$ of Sec.~\ref{sec:defs}.C is therefore diagonal, with $H_{{\rm ang},0}$ on the $++$ structure and on the
parity-odd structure of even $J$, and $H_{{\rm ang},4}$ on the $+-$ structure and on the parity-odd structure of odd $J$, each times a
positive factor. Hence (Pos$_\Omega$) on $\mathcal O$ is $H_{{\rm ang},0}\ge0$ for even $J$ and $H_{{\rm ang},4}\ge0$ for $J\ge4$.
\emph{Grade.} Eq.~\eqref{eq:fp-target} is [D given the ladder and first-moment identities of Sec.~\ref{sec:angular} as operator
identities]. Its algebra was checked symbolically, and the Einstein calibration of $\mathcal B$ at $k=2$ in exact arithmetic; the
$k=3$ prefactor and the flip cancellation are [D]. Only the sign of $\mathcal B$ is used.

\emph{Congruence lift of T0.} Define $\Omega^{\rm hyb}_{\mathcal O}$ and $N^{\rm phys}_{s,\mathcal O}$ by
Eq.~\eqref{eq:fp-target} with $H_{{\rm ang},b}$ replaced by T0's hybrid action $H_{{\rm hyb},b}$ and by $N_{s,b}$. Since T0 gives
$H_{{\rm hyb},b}\ge0.1999\,N_{s,b}$,
\begin{equation}
|H_{{\rm ang},b}-H_{{\rm hyb},b}|\le(0.1999-\epsilon)\,N_{s,b}\ \ (b=0,4)\quad\Longrightarrow\quad
\Omega^{EE}_{\mathcal O}\ge\epsilon\,N^{\rm phys}_{s,\mathcal O}\ge0
\label{eq:fp-lift}
\end{equation}
[P]; $\mathcal B$ cancels. The comparison is meaningful because the leading-Regge limit of $H_{{\rm ang},b}$ is T0's per-state
action with relative factor exactly 1 (Sec.~\ref{sec:fullpos}.B). The certificates of Sec.~\ref{sec:fullpos}.C prove the sign of
$H_{{\rm ang},b}$ directly, without the comparison; the comparison is used only as evidence.

\subsection{The exact finite-mass action}
\emph{Kernels.} Let $\delta=\tfrac12+i\nu$ and let $K_\delta(\varphi)$ be the celestial factor of the dual $EE$ structure (gauge
$x_0=\infty$) integrated over the mean direction of the two detectors at relative angle $\varphi$. Exact integrations by parts of the
insertions $\int(\cos t-c)^j|\cos t-c|^{-\delta}dt$ give the $k=2$ kernels of $\Omega^{EE}$ at finite mass, with no large-$\Delta$,
large-$J$ or large-$\nu$ approximation:
\begin{equation}
\mathcal Z=12K_\delta-2\sin\varphi\,K_\delta',\qquad
\mathcal Y=\Big(6-\frac{\nu^2+\frac14}{2}\Big)(1+\cos\varphi)K_\delta-7\sin\varphi\,K_\delta'
\label{eq:fp-kernels}
\end{equation}
[D; the Fourier algebra below P]. The Fourier coefficients of $K_\delta$ obey the continuous dual Hahn (CDH) recurrence, so they are
proportional to $f_m(\nu)$ defined by $f_0=0$, $f_1=1$, $(m+1)^2f_{m+1}=-(s-2m^2)f_m-(m-1)^2f_{m-1}$, $s=\nu^2+\tfrac14$. Those of
$\mathcal Z$, $\mathcal Y$ are
\begin{equation}
z_m=12f_m-(m-1)f_{m-1}+(m+1)f_{m+1},\qquad
y_m=\tfrac{12-m^2}2f_m+\tfrac{13-7m}2f_{m-1}+\tfrac{13+7m}2f_{m+1}+\tfrac{(m-2)^2}4f_{m-2}+\tfrac{(m+2)^2}4f_{m+2},
\label{eq:fp-zy}
\end{equation}
which are the coefficients of the first-moment construction of Sec.~\ref{sec:angular}: the angular-route coefficients are exactly the
$k=2$ kernels of $\Omega^{EE}$. For $k=3$ the $d=1$ CDH polynomials enter, $q_0=1$,
$(n+4)^2q_{n+1}=-(s+4-2(n+2)^2)q_n-n^2q_{n-1}$.

\emph{CDH generating function} [D; DV]. With $f_m=Q^{(0)}_{m-1}$ and $q_n=Q^{(1)}_n$,
\begin{equation}
\sum_{n\ge0}\Big(\frac{(2d+2)_n}{n!}\Big)^2r^nQ^{(d)}_n(\nu^2)=\frac{P_{-1/2+i\nu}\big(\frac{1+r}{1-r}\big)}{(1-r)^{2d+2}},\qquad d=0,1.
\label{eq:fp-cdh}
\end{equation}
The proof turns the recurrence into an ODE that reduces to the Legendre equation; the identity, and its $k=2$ first-moment form
(identity I2 below), hold to $2.4\times10^{-40}$ and $6.5\times10^{-40}$ relative at twelve $(r,\nu)$ points.

\emph{Moments.} With the $k=2$ continuum weight $\chi$ and the cutoff $\theta$ of Sec.~\ref{sec:defs},
\begin{equation}
g(\nu)=-\frac{9216}{\Lambda}\,\frac{4\nu^2+1}{(4\nu^2+9)(4\nu^2+49)}\,\chi(\nu)\theta(\nu),\quad
Z_k=\int_0^\infty\!g\,z_{k+2}\,d\nu,\quad Y_k=\int_0^\infty\!g\,y_{k+2}\,d\nu,\quad Q_k=\int_0^1\!\psi_3(p)\,q_k(\Lambda^2p^2)\,dp,
\label{eq:fp-moments}
\end{equation}
with $\psi_3$ the NF02 weight, so that the $Q_k$ are exact rationals.

\emph{The action.} With $V_{0,k}=4k/(k+3)$ and $V_{4,k}=4+12/(k+1)$,
\begin{equation}
H_{{\rm ang},b}=(1-r)^3\sum_{k=0}^{n}D_k^2\,U_{b,k}\Big[Z_k\Big(\frac1{k+3}-\frac r{k+1}\Big)+\frac{2Y_k+V_{b,k}Z_k+\sigma_b\Lambda^2Q_k}{T}\Big]
=:H^{(2)}_{{\rm ang},b}+H^{(3)}_{{\rm ang},b},
\label{eq:fp-H}
\end{equation}
where $H^{(3)}$ is the $Q_k$ part. At fixed $J$ this is a rational function of $T$, with no approximation.

\emph{Single-sequence form} [P; checked in exact rational arithmetic]. Put $L_k=(Z_k-\frac k{k+3}Z_{k-1})/(k+3)$ ($L_0=Z_0/3$) and
$P_k=D_k^2L_k$. Summation by parts in $k$ gives $H^{(2)}=(1-r)^3\sum_kD_k^2U_{b,k}[L_k+(2Y_k-u_{k-1}L_k)/T]$. For any sequence $f_m$,
Eq.~\eqref{eq:fp-zy} implies $2Y_k=\frac{(k+3)^2}2L_k+\frac{(k+4)^2}2L_{k+1}$, which eliminates $Y_k$. Hence
\begin{equation}
H^{(2)}_{{\rm ang},b}=(1-r)^3\sum_{k=0}^{n+1}\mathcal V_{b,k}P_k,\qquad
\mathcal V_{b,k}=U_{b,k}\Big[1+\frac{(k+3)^2/2-u_{k-1}}{T}\Big]+\frac{k^2}{2T}U_{b,k-1}\ \ge0,
\label{eq:fp-single}
\end{equation}
where at $k=0$ the value $u_{-1}=(b-1)b$ must be used ($12$ for $b=4$), i.e.\ the formula for $u_j$ continued to $j=-1$; with
$u_{-1}=0$ the identity fails for $b=4$, by about $10\%$. $\mathcal V_{b,k}\ge0$ because $u_{k-1}\le C<T$ for $k\le n+1$. We checked
Eq.~\eqref{eq:fp-single} exactly on 40 random states with random sequences $f_m$ (with $u_{-1}=0$ all 21 cases with $b=4$ fail), and
with the certified moments it agrees with Eq.~\eqref{eq:fp-H} to $10^{-536}$ at $J=10$, $40$, $100$. In generating-function form,
with $\mathcal L$ the linear map $\mathcal L(z^k):=U_{b,k}$ and $F(z):=\sum_kP_kz^k$,
\begin{equation}
H^{(2)}_{{\rm ang},b}=(1-r)^3\,\mathcal V_b(F),\qquad \mathcal V_b:=\mathcal L\circ\Big(1+\frac{\mathcal D_b}T\Big),\qquad
\mathcal D_b=\frac{z(1-z)}2\partial_z^2+\Big[\Big(\frac72-2b\Big)z+\frac12\Big]\partial_z+\frac92-b(b-1),
\label{eq:fp-V}
\end{equation}
$\mathcal V_b(z^k)=\mathcal V_{b,k}$ [P]. $\mathcal D_b$ is the first-moment operator in generating-function form; it is unrelated to
the spectral defect $\mathcal D$ of Sec.~\ref{sec:angular}, and its constant term is the $u_{-1}$ value above. The generating function
itself is $F(z)=\frac2{\Lambda(1-z)^3}\int E_\theta(\nu)P_{-1/2+i\nu}\big(\frac{1+z}{1-z}\big)d\nu$, with $E_\theta$ the shifted
density of the realised weight [D].

\emph{Leading-Regge limit, unit relative normalisation.} As $T\to\infty$ at fixed $r$, $U_{b,k}\to r^k$ and every $1/T$ term drops.
Identity I2, $(1-r)^3\sum_kD_k^2r^kz_{k+2}(\frac1{k+3}-\frac r{k+1})=-\frac{(4\nu^2+9)(4\nu^2+49)}{4608(4\nu^2+1)}[R_-P_{-1/2+i(\nu+i)}(\eta)
+R_+P_{-1/2+i(\nu-i)}(\eta)]$ with $\eta=(1+r)/(1-r)$, and Eq.~\eqref{eq:fp-cdh} at $d=1$ give
\begin{equation}
H_{{\rm ang},b}\ \to\ \frac2\Lambda\int_0^\infty\!\chi\theta\,\big[R_-P_{-1/2+i(\nu+i)}(\eta)+R_+P_{-1/2+i(\nu-i)}(\eta)\big]d\nu
+\sigma_b\,u\!\int_0^1\!\psi_3(p)P_{-1/2+i\Lambda p}(\eta)\,dp ,
\label{eq:fp-LR}
\end{equation}
$u=\Lambda^2/m^2$. The first term is $2/\Lambda$ times the realised leading-Regge image of the $k=2$ weight; by the de-shift identity it
equals $2\Omega_2$ up to $2\varepsilon_WP_{-1/2}(\eta)/\Lambda\le5.4\times10^{-9}$ (certified $\varepsilon_W$ of Sec.~\ref{sec:T0}).
So the limit is $F_{\rm LR}(\beta,u,\sigma_b)=2\Omega_2+u\sigma_b\Omega_3$ of Sec.~\ref{sec:largespin}, with relative normalisation
exactly 1; the factor 2 is $9216/4608$ [D; DV]. Checks: at $J=0$, $T\to\infty$, $H_{{\rm ang},0}=Z_0/3=20.7718930153=2\int_0^1\psi_2$
[certified]. At fixed $r$ the finite-$T$ deviation scales like $1/T$, e.g.\ $-1.1\times10^{-4}$ and $-2.8\times10^{-5}$ relative at
$\hat b=2$ for $J=20$ and $J=40$ ($T$ four times larger) [DV]. The difference between $\eta$ and $\cosh\beta=\eta+1/(2m^2)$ is of the
same order.

\subsection{Certified coverage at \texorpdfstring{$\Delta_g=1000$}{Delta g = 1000}}
\emph{Certified moments} [certified]. $Z_k$, $Y_k$ are computed in Arb ball arithmetic in closed form, with no quadrature of the
Gaussian part of $\chi$. That part is reduced by exact polynomial division to exact Gaussian moments and Faddeeva functions. The
logarithmic part is reduced to an exact integration-by-parts recurrence and three one-dimensional integrals done with Arb's rigorous
integrator. The cutoff, whose relative effect on the $k$-th moment is about $10^{-24}\times1.03^{2k}$ and therefore not negligible,
is treated by Gauss--Legendre quadrature near $\nu=\Lambda+100$ with the Bernstein-ellipse remainder bound, and by Gaussian tail bounds
elsewhere. At $\Delta_g=1000$ the moments are certified for $k\le1000$ (7000-bit run, 20\,min), with relative radius $\le10^{-500}$ for
$k\le500$ and $\le10^{-216}$ for $k\le1000$; at $\Delta_g=300$ for $k\le300$. \emph{Rigour patch.} In the first runs the logarithm inside
Arb's integrator was evaluated without the flag that enforces the holomorphy check on which its error bound rests. Every moment was
recomputed with the flag, and every recomputed ball lies inside the ball used by the certificates: 2002 of 2002 ($k\le1000$), and 602 of 602
for $k\le300$ at both $\Delta_g=1000$ and $300$. So all certificates stand as run.

\emph{Fixed spin, every mass.} With $\varpi:=\Lambda^2/T$, $H_{{\rm ang},b}\prod_{j<n}(T-u_j)/[(1-r)^3T^n]$ is a polynomial
$\mathcal P_{J,b}(\varpi)$ of degree $n+1$, and the prefactor is positive on heavy states. The heavy range of spin $J$,
$m\ge m_{\rm gap}(J)$, is $0\le\varpi\le\Lambda^2/[(J+\Delta_g-1)(J+\Delta_g-2)]$, including $T\to\infty$. $\mathcal P_{J,b}>0$ there is
proved by positivity of its Bernstein coefficients (Arb, from the moment balls and the exact $Q_k$), with de Casteljau bisection where
needed. Every channel with $J\le766$ passes: 1147 polynomials, 356 for $J\le239$ with the $k\le300$ moments and 791 for
$240\le J\le766$ with the $k\le1000$ moments. There is no failure; the runs were stopped at $J=766$ for time. With the $k\le300$ moments
the first failure is at $J=240$, from moment precision, not a sign. The Bernstein margin is smallest at the twist edge and falls with
$J$ ($\mathcal P(\varpi_{\max})\approx0.52$, $2.2\times10^{-3}$, $2\times10^{-6}$ at $J=4$, $40$, $238$), through the normalising factor
$\prod_j(1-u_j\varpi/\Lambda^2)$, not through $H_{{\rm ang},b}$. At $\Delta_g=300$ every channel with $J\le285$ passes (426
polynomials; the count includes the $b=4$ channel of $J=286$); the first failure, at $J=286$, $b=0$, is again from moment precision.

\emph{All spins at bounded impact parameter.} In the variables $x=\hat b^2/4=\Lambda^2(J+\frac12)^2/m^2$ and $v=(J+\frac12)^{-2}$,
exactly, $C=1/v-\frac14$, $1/T=xv/E$ with $E=\Lambda^2+x-xv/4$, and
\begin{equation}
U_{b,k}=\frac{x^kA_{b,k}(v)}{\Lambda^{2k}\Pi_k},\qquad A_{b,k}=\prod_{j<k}\big[1-(b+j+\tfrac12)^2v\big],\qquad
\Pi_k=\prod_{j<k}\Big[1+\frac x{\Lambda^2}\big(1-(b+j+\tfrac12)^2v\big)\Big].
\label{eq:fp-xv}
\end{equation}
The terms $k\le K$ of Eq.~\eqref{eq:fp-H}, multiplied by $E\Pi_K>0$, form an exact bivariate polynomial $N(x,v)$ of degree $K+1$ in each
variable; no denominator is expanded. For $J\ge J_0>K+b$ every such term is a genuine ladder term, and termination is automatic
($A_{b,k}=0$ for $k>n$). The terms $K<k\le1000$ are bounded with the certified moments, $0\le A_{b,k}\le1$, $\Pi_k\ge1$, and those with
$k>1000$ by the all-$k$ bounds $|f_m(\nu)|\le3^{m-1}e^{(\pi/\sqrt3)(\nu+\frac12)}$ and $|q_n|\le3^ne^{(\pi/\sqrt3)\sqrt{\nu^2+17/4}}$
integrated in closed form. Tensor Bernstein coefficients of $N$, less the tail bound, are positive on $[0,x_{\max}]\times[0,(J_0+\frac12)^{-2}]$
for both channels, for the three boxes of Table~\ref{tab:fpcover} [certified]. The method stops at the harmonic truncation $K<J_0-b$:
with $|Z_k|\lesssim10^{673}$ the tail $x^K|Z_K|/\Lambda^{2K}$ is small only for $\hat b\lesssim0.7J_0$. (A box $J_0=651$,
$\hat b\le490$ failed because of the crude $k>1000$ bound, not because of $N$.)

\begin{table}[!htbp]
\caption{\label{tab:fpcover}Positivity of the finite-mass action, $H_{{\rm ang},b}>0$ in both channels, at $\Delta_g=1000$ unless stated;
$K$ is the number of harmonics kept, TB the certified tail bound, and $N(0,0)=2.07\times10^7$ in the bivariate rows.
$\hat b_{\max}(J)=(2J+1)\Lambda/\sqrt{(\Delta_g-2)(\Delta_g-1+2J)}$ is the impact parameter at the twist edge ($\approx965$ at $J=767$,
$\approx44.7\sqrt J$ at large $J$); $\ell=-\ln(1-r)$.}
\footnotesize
\begin{ruledtabular}
\begin{tabular}{L{4.6cm}L{9.0cm}L{2.6cm}}
region & method & grade\\\hline
$J\le766$, every $m\ge m_{\rm gap}(J)$ & per-spin Bernstein certificate in $\varpi=\Lambda^2/T$; 1147 polynomials, no failure & [certified]\\
$J\ge240$, $\hat b\le100$ & bivariate $(x,v)$, $K=230$, TB $=9.3\times10^{-122}$; min.\ Bernstein coefficient $1.56\times10^3$ ($b=0$), $1.55\times10^3$ ($b=4$) & [certified]\\
$J\ge420$, $\hat b\le282.8$ & bivariate, $K=410$, TB $=7.3\times10^{-54}$; $1.56\times10^3$, $1.53\times10^3$ & [certified]\\
$J\ge651$, $\hat b\le400$ ($\ell\le0.039$) & bivariate, $K=640$, TB $=9.8\times10^{-90}$; $2.93\times10^3$, $2.89\times10^3$ & [certified]\\
$\Delta_g=300$: $J\le285$, every mass & per-spin Bernstein certificate, 426 polynomials & [certified]\\
$\ell\ge30.16$ ($\hat b\gtrsim7.1\times10^9$), every $J$ & contour argument of Sec.~\ref{sec:fullpos}.D, uniform transform bound & [D-sketch]\\
$20.27\le\ell<30.16$ & the same with the transform profile, itself open & [I]\\
$J\ge767$, $400<\hat b\le\hat b_{\max}(J)$, $\ell<30.16$ & open; numerical values at six states with $\ell\le0.29$ (Sec.~\ref{sec:fullpos}.E) & [I]; values [DV]\\
\end{tabular}
\end{ruledtabular}
\end{table}

\emph{Comparison with the hybrid} [DV]. On 69 state--channel pairs ($J\in\{0,2,4,5,6,8,10,12,16,20,30,40,60\}$, both channels,
twists 1000, 1300, 2000), with $H_{{\rm ang},b}$ from the certified moments and T0's hybrid in floating point,
$\max|H_{{\rm ang},b}-H_{{\rm hyb},b}|/N_{s,b}=4.07\times10^{-3}$ ($J=16$, $+-$, twist edge); it is $\le4.7\times10^{-5}$ for $++$ and
$\le4.2\times10^{-5}$ for $+-$ at twist $\ge1300$. Sampling all 37 channels with $J\le26$ at 24 masses each gives
$\max|H_{{\rm ang},b}/H_{{\rm hyb},b}-1|=4.71\times10^{-3}$ ($J=17$, $+-$, twist edge) and $\min H_{{\rm ang},b}/N_{s,b}=0.2009$ ($J=5$,
$+-$). A sharper statement, $|H_{{\rm ang},b}/H_{{\rm hyb},b}-1|\le0.01$ for $J\le26$ at every mass, was obtained with a separate
certificate in the variable $u$ that keeps the exact Weyl shift; it has not been reproduced, and we do not use it. So where it was measured, the
finite-mass correction uses at most $2\%$ of the T0 margin $0.1999$.

\subsection{Large impact parameter}
As $r\to1$ the ladder sum is controlled through singular powers rather than through harmonics. Near $z=1$ the Legendre function in
$F(z)$ is a combination of $((1+z)/(1-z))^{-1/2\pm i\nu}$, so $F$ is a superposition over $\nu$ of $(1-z)^{-\alpha}$, $\alpha=\frac52+i\nu$,
with density $\mathsf f(\nu)\mathsf B(\nu)$, $\mathsf B(\nu)=\Gamma(2i\nu)/\Gamma(\frac12+i\nu)^2$,
$\mathsf f(\nu)=\chi_\theta R_-(\nu-i)+\chi_\theta R_+(\nu+i)$ (on the real axis $\mathsf f$ reproduces $\psi_2$) [D]. On each power,
Eq.~\eqref{eq:fp-V} gives exactly
\begin{equation}
\mathcal V_b\big((1-z)^{-\alpha}\big)=\Big(1+\frac{B_b}T\Big)S_b(\alpha)+\frac{\alpha A_b}T\big[S_b(\alpha+1)-S_b(\alpha)\big],\qquad
S_b(\alpha)=\sum_kU_{b,k}\frac{(\alpha)_k}{k!},
\label{eq:fp-Salpha}
\end{equation}
$A_b=\frac{9+\alpha}2-2b$, $B_b=\frac{9+\alpha}2-b(b-1)$ [P; DV to $6\times10^{-41}$]. In the leading-Regge limit
$S_b(\alpha)\to w^{-\alpha}$. The deviation $\mathsf D_\alpha:=w^\alpha\mathcal V_b((1-z)^{-\alpha})-1$ enters the leading connection
term
\begin{equation}
\mathcal E_2=\frac4\Lambda\sqrt w\;{\rm Re}\!\int\!\mathsf f(\nu)\mathsf B(\nu)\,e^{i\nu\ell}\,\mathsf D_{5/2+i\nu}\,d\nu .
\label{eq:fp-E2}
\end{equation}
For real $0\le\alpha\le4$ and $r\ge\frac12$, $|\mathsf D_\alpha|\le94/a^2$ on every heavy state (a uniform singular-power bound) [D;
tested on 40 states, worst $14.8/a^2$]. For complex $\alpha$ there is no bound polynomial in $|\nu|$, since
$\mathcal V_b((1-z)^{-\alpha})$ has exact degree $n+1$ in $\alpha$ [P]. An absolute-value $\nu$-integration of Eq.~\eqref{eq:fp-E2} cannot
work at moderate $\ell$ [D]. The estimate is therefore signed. The $\nu$-contour is deformed to
$0\to\frac34i\to\frac14+\frac34i\to\frac14+\frac94i\to\infty+\frac94i$. On the imaginary segment the integrand defines a density
$W_{\rm loc}(y)=-2\mathsf B(iy)\,{\rm Im}\,\mathsf f(iy)$, $0\le y\le\frac34$; with the local bound $|\mathsf D(iy)|\le100/a^2$ it
contributes at most $(100/a^2)\mathcal M(\ell)$, where
\begin{equation}
\mathcal M(\ell)=\sqrt w\Big[\mathsf A+\frac2\Lambda\int_0^{3/4}W_{\rm loc}(y)\,e^{-\ell y}\,dy\Big],\qquad \mathsf A=E(0)/\Lambda\ge0,
\label{eq:fp-M}
\end{equation}
is the positive main term. $\mathsf A\ge0$ is the endpoint condition $E(0)\ge\psi(0)$ of Sec.~\ref{sec:T0}; without the repair it would
fail. The detour and the horizontal leg are damped by $e^{-3\ell/4}$ and $e^{-9\ell/4}$. On them the argument uses a finite-state
beta-kernel bound $|K_{n,a,b}(q)|\le2|1+q|^{-a}$, proved for $n\ge1024a^2$ and stated for $n<1024a^2$ when $\ell\ge40$, and a transform
bound on the leg, $|\mathsf D(t+\frac94i)|\le200(1+t)^{5/4}e^{a\,g(t/a)}$ for $t\ge0$, whose supremum on the band is $<5\times10^{12}$
[D-sketch].

\emph{Certified contour inputs} [certified, Arb]. $\sup|\mathsf f(t+\frac34i)|\le3.924$ for $0\le t\le\frac14$; on the vertical leg,
$\sup_y|\mathsf f(\frac14+iy)|e^{-17(y-3/4)}\le227$ (the twelve values of $\sup|\mathsf f|$ on $\frac18$-intervals are also enclosed,
rising from $12.6$ to $2.69\times10^{13}$); $\int_{1/4}^\infty|\mathsf f(t+\frac94i)|dt\le9.0084\times10^{13}$ and
$\int_{1/4}^\infty|\mathsf f\mathsf B|(t+\frac94i)\,dt\le5.3018\times10^{11}$; $W_{\rm loc}=(\frac12-y)^2V(y)$ with $V>0$ on $[0,\frac34]$, so
$W_{\rm loc}\ge0$ with a double zero at $y=\frac12$ (the sign changes of $\mathsf B$ and $R_+$ there cancel), $W_{\rm loc}(0)=0.0604088$,
and $\int_0^{3/4}W_{\rm loc}e^{-\ell y}dy\ge1.945\times10^{-3}$ at $\ell=28$; $|\mathsf B|\le0.1034$ on the detour and $\le0.0171$ on the
leg. With these, $|\mathcal E_2|/\mathcal M(\ell)\le\rho(\ell):=100/a^2+2X(\ell)/I_W(\ell)$, where $X(\ell)$ collects the three contour
pieces and $I_W$ is the integral in Eq.~\eqref{eq:fp-M}. $\rho$ is evaluated rigorously on a cell cover of $28\le\ell\le400$, and
beyond $\ell=400$ every term decays.

\emph{Thresholds.} With the uniform transform bound $5\times10^{12}$, $\rho<0.01$ for $\ell\ge30.16$ and $\rho<0.1$ for $\ell\ge29.12$,
while $\rho(28)\le1.19$. With the profile, $\rho<0.01$ for $\ell\ge20.27$ and $\rho(28)\le1.004\times10^{-4}$, the floor $100/a^2$. About
99\% of $\rho$ comes from the horizontal leg, where the transform bound, not $\mathsf f$, is now the limiting input. Assembled with the
bounds on the remaining connection terms, a localisation lemma $|\mathcal L(G)|\le50K(4w)^p$ for $|G(z)|\le K|1-z|^p$ ($1\le p\le597$),
and on the $k=3$ part, $|H^{(3)}|<4\times10^5w^{3/2}$, all of which decrease with $\ell$, this gives $H_{{\rm ang},b}>0$ in both channels
for every spin and mass with $\ell\ge30.16$, i.e.\ $\hat b\gtrsim7.1\times10^9$. At the twist edge that is $J\gtrsim2.5\times10^{16}$.
With the profile it would start at $\ell\ge20.27$ ($\hat b\gtrsim5.0\times10^7$; $J\gtrsim1.3\times10^{12}$ at the edge).
\emph{Grade.} $\ell\ge30.16$: [D-sketch]. The contour inputs are [certified]; the transform bound, the local bound $100/a^2$, the
kernel bounds, the connection terms and the $k=3$ bound are written-out but unchecked derivations, and the kernel bound for
$n<1024a^2$ is stated for $\ell\ge40$ only. $20.27\le\ell<30.16$: [I], because it needs the transform profile below its supremum,
the complex-$\alpha$ analogue of the singular-power bound, which is not proved. (An earlier claim of positivity for $\ell\ge28$
rests on the profile and is not supported by the uniform bound.)

\subsection{Why it is hard, and what remains}
\emph{Exponentially large moments.} $|Z_k|$ grows like $(\Lambda^k/k!)^2$ up to $k\approx\Lambda/2$ and then stays on an oscillating
plateau, $|Z_k|\approx10^{672.5}$ for $550\le k\le1000$ ($10^{198.5}$ for $160\le k\le300$ at $\Delta_g=300$) [certified values], against
$e^{(\pi/2)\Lambda}=10^{681}$ ($10^{204}$). The reason is the CDH polynomials: at fixed $\nu$, $f_m$ grows to about $e^{(\pi/2)\nu}$ for
$m<\nu/2$ (since $\int_0^1{\rm arcosh}(1/s)\,ds=\pi/2$) and then oscillates with amplitude $e^{(\pi/2)\nu}m^{-3/2}$ [D; DV]. The ladder sum of
size $10^{-5}$--$10^{1}$ is therefore a cancellation among terms of size up to $10^{685}$. Consequences: the per-spin certificate needs
moments to about $J/2+300$ digits; coefficient envelopes of the type $|P_k-P^*_k|\le P^*_k/4$, with
$P^*_k=((5/2)_k/k!)(\mathsf A+\mathsf B'/\ln(k+2))$, hold for no $k\le1000$ and can hold only for $\ln k\gtrsim\Lambda$ [certified values; D];
the absolute comparison constant $\mathcal K_b(r_0)$ of a bound $|H_{\rm ang}-H_{\rm LR}|\le\mathcal K_b(r_0)/T$ grows like $e^{2\hat b}$
($\ge1.8\times10^7$, $7\times10^{11}$, $1.8\times10^{27}$, $3.9\times10^{84}$ at $r_0=10^{-6}$, $10^{-4}$, $10^{-3}$, $10^{-2}$), so it is
useless beyond $\hat b\approx20$; and a signed remainder in $r$ carries the same cancellation. Per frequency, in contrast, the finite-mass
kernel is of polynomial size and within $O(\nu^2/T)$ of the leading-Regge kernel [DV]: the exponential sizes live only in the moment
basis.

\emph{$\mathcal L$ is not a positive average.} With $c:=U_{b,1}=(C-u_0)/(T-u_0)$ (so $c=r$ for $b=0$), $\mathcal L(z-c)=0$ and
\begin{equation}
\mathcal L\big((z-c)^2\big)=U_{b,2}-c^2=-\frac{2(b+1)\,c(1-c)}{T-u_1}<0
\label{eq:fp-notpos}
\end{equation}
[P; checked exactly on 320 states]. So $\{U_{b,k}\}$ is not the moment sequence of a positive measure, and positivity of a generating
function on $[0,1]$ does not transfer to its ladder sum: no monotonicity argument can carry T0's floor to finite mass.

\emph{The fixed-twist limit falls short.} As $J\to\infty$ at fixed $a$ and fixed $\xi=k(1-r)$, $U_{b,k}\to e^{-2a\,{\rm artanh}(\xi/2a)}$, and
the finite-mass action of the leading singular power, relative to its leading-Regge value, tends to
\begin{equation}
\mathfrak N_a=\frac1{\Gamma(5/2)}\int_0^{2a}\xi^{3/2}e^{-2a\,{\rm artanh}(\xi/2a)}d\xi,\qquad 1-\frac{105}{32a^2}\le\mathfrak N_a<1
\label{eq:fp-Na}
\end{equation}
[D; DV: $\mathfrak N_{998}=0.999996705605$ against the lower bound $0.999996705586$]. The deficit is small, $3.3\times10^{-6}$ at the
twist edge, but it does not vanish as $J\to\infty$: the finite-mass correction has no favourable sign and does not decay at fixed twist,
so uniform control needs explicit $1/a^2$ bounds such as the one quoted after Eq.~\eqref{eq:fp-E2}. For complex
$\alpha=\frac52+i\nu$ the analogue of $\mathfrak N_a$ has a turning point at $|\nu|=a$ and grows like $e^{a\,g(\nu/a)}$ beyond it [D].

\emph{First-order profile.} With $d:=1-c$, $X:=T-u_0$ (so $dX=m^2$) and $S_{b,k}:=\sum_{j<k}(u_j-u_0)=k(k-1)(k+3b+1)/3$,
\begin{equation}
\big|U_{b,k}-c^k+(d/X)\,c^{k-1}S_{b,k}\big|\le6\,c^{k-2}S_{b,k}^2/X^2\qquad\text{for all }k\ge0,
\label{eq:fp-ladder}
\end{equation}
including $k>n$, where $U_{b,k}=0$ [D; checked exactly]. We checked Eq.~\eqref{eq:fp-ladder} in exact rational arithmetic on 700 heavy
states, $b=0$, $4$, at gaps $\Delta_g=6$, $30$, $300$, $1000$, with spins up to $10^6$ and twists up to $10^7$: 173\,009 pairs $(\text{state},k)$,
17\,350 of them on terminated ladders, every $k\le n+30$ on the 380 states with $n\le1500$. There was no failure, and the largest ratio of the
two sides was $0.082$. Summed against $F=\sum_kp_kz^k$ it gives
\begin{equation}
\mathcal L(F)=F(c)+\frac{\mathcal A_bF(c)}{T-u_0}+\mathcal R,\qquad
\mathcal A_bF=-c(1-c)\Big[(b+1)F''+\frac c3F'''\Big],\qquad
|\mathcal R|\le\frac6{(T-u_0)^2}\sum_k|p_k|\,c^{k-2}S_{b,k}^2 .
\label{eq:fp-profile}
\end{equation}
To first order the ladder evaluates $F$ at $c$, with an explicit correction that is negative on functions with positive derivatives,
consistent with Eqs.~\eqref{eq:fp-notpos} and~\eqref{eq:fp-Na}. On a singular power the relative correction is
$-c\,\alpha(\alpha+1)[(b+1)+c(\alpha+2)/(3(1-c))]/m^2$. Near the twist edge $m^2(1-c)\to4a^2$ as $J\to\infty$, so the correction is of order
$|\nu|^3/(12a^2)$ and becomes of order one at $|\nu|\approx230$, well inside the support $\nu\lesssim\Lambda$ of the weight. Away from the
far tail, where $e^{i\nu\ell}$ is used through the contour, the first-order profile is therefore not a small-parameter expansion of
$H_{\rm ang}$, and its remainder has to be controlled with its cancellations.

\emph{What remains open.} At $\Delta_g=1000$ positivity of the finite-mass action is open exactly on
\begin{equation}
\mathcal R_{\rm open}=\Big\{J\ge767,\ \ 400<\hat b\le\hat b_{\max}(J)\approx44.7\sqrt J,\ \ \ell<30.16\Big\},
\label{eq:fp-open}
\end{equation}
that is $m_{\rm gap}(J)\le m<(2J+1)\Lambda/400\approx4.99(J+\frac12)$ at large spin, including $r\to1$ at fixed twist. For
$J\lesssim2.5\times10^{16}$ the constraint $\ell<30.16$ is automatic. With the transform profile, $30.16$ would become $20.27$. The region is infinite.
Every finite range of $J$ can be closed at polynomial cost, with moments to $k\sim J$ and bivariate boxes with $K\approx J_0$, but the
edge needs infinitely many boxes, and its limit, $J\to\infty$ at fixed twist, is reached by no finite computation. Two analytic routes
avoid the moments: a per-frequency bound on the finite-mass kernel minus the hybrid kernel, $O(\nu^2/T)$ with no exponential sizes,
or a Mellin representation of the ladder sum with its boundary term at $r=1$ resolved. Since $\hat b>400\ge0.37211\Lambda$ there,
Theorem~L$'$ of Sec.~\ref{sec:largespin} covers $\mathcal R_{\rm open}$ for $J\le10^6$ at its conditional grade (on premises unrelated to
this section).

\emph{Numerical evidence in $\mathcal R_{\rm open}$} [DV]. Exact ladder sums from the certified moments, at $b=0$, $4$ and twists
$a=998$--$10^5$ with $n\le1000$, reach $\ell\le0.29$. Of the 18 states scanned, six lie in $\mathcal R_{\rm open}$ ($J=1000$, $1004$;
$a=998$, $1500$, $3000$; $\ell=0.065$--$0.29$). There $H_{{\rm ang},b}/N_{s,b}\ge0.918$ ($b=4$, $a=998$; $\approx1.0001$ for $b=0$), against
the T0 floor $0.1999$, and on all 18 states $H_{{\rm ang},b}$ agrees with the hybrid to $\le1.3\times10^{-4}$ relative; the smallest value in
the scan is $0.893$ ($b=4$, $a=998$, $J=445$, $\ell=0.10$, inside the certified region). Together with the hybrid comparison of
Sec.~\ref{sec:fullpos}.C ($|H_{{\rm ang},b}-H_{{\rm hyb},b}|/N_{s,b}\le4.07\times10^{-3}$ on 69 states with $J\le60$), this shows the
finite-mass correction far below the margin wherever it has been computed. No state with $0.29<\ell<30.16$ in $\mathcal R_{\rm open}$
has been evaluated.

\emph{For assumption (iii).} At $\Delta_g=1000$ the finite-mass action of $\Omega^{EE}$ is certified positive on every heavy state
outside $\mathcal R_{\rm open}$ with $\ell<30.16$, that is everywhere except near the twist edge at spins $J\ge767$ and impact parameters
$400<\hat b$ ($\ell\ge30.16$ is covered at grade [D-sketch]). In that large-spin, near-edge region, the states computed (all with
$\ell\le0.29$) keep a margin far from zero. The identification of this action with $\Omega^{EE}$ is [D given the operator identities of
Sec.~\ref{sec:angular}].

\emph{Code.} \texttt{code/fullpos/} (moments, per-spin and bivariate certificates, comparisons, scan) and \texttt{code/contour/} (contour
inputs). The certified moments, the per-spin and bivariate certificates and the contour inputs have not yet been reproduced by separate
code.
